% ════════════════════════════════════════════
% 模型检验（06.模型检验.tex）
% ════════════════════════════════════════════
\section{模型检验}

为验证最终模型的有效性与稳健性，本节按两条线索逐项检查：先报告拟合优度、检验整体显著性并做稳定性与基线比较，再计算训练集与测试集误差、判断过拟合并诊断残差，最后用七种设定复核核心效应。检验不替换最终模型，只观察同一组结论在不同口径与不同设定下是否保持一致。

\subsection{模型有效性与稳定性}

最终模型 M2 含 26 个参数、500 条观测，样本内 $R^2$ 为 0.2068、调整 $R^2$ 为 0.1650，残差标准误为 1.0765。全部解释变量合计解释约两成曝光方差，其余部分来自未被观测的推荐机制与传播随机性；调整 $R^2$ 与 $R^2$ 相差不大，说明 26 个参数没有靠数量堆出拟合。M2 的整体 $F$ 统计量为 4.9443、$p<0.001$，拒绝“全部斜率同时为零”的原假设。分项看，付费、Status、Video 与主题 2、主题 3 的系数均显著，Link、时段与星期的系数不显著。整体显著只说明模型整体有效，不代表每个参数都可识别，粉丝规模因共线被剔除的依据见表\ref{tab:q3_diag}。

十折交叉验证重复 5 次、共 50 个测试折。M2 的 CV $R^2$ 为 0.1010，高于 M1 的 0.0972、M3 的 0.0935 与 M4 的 0.0972，更高于均值基线 M0 的 $-0.0212$；CV RMSE 为 1.1007，是五种设定中最低。样本外表现优于全部对照设定，说明剔除粉丝规模、保留月份固定效应的简化没有牺牲预测能力。

\subsection{误差分析与残差诊断}

训练集误差在全样本上计算，测试集误差取 50 个测试折的平均。M2 的训练集 RMSE 为 1.0482、MAE 为 0.7918；测试集 RMSE 为 1.1007、MAE 为 0.8328，两项绝对误差都只上升约 5\%，说明模型没有出现严重的过拟合。$R^2$ 由 0.2068 降至 0.1010，降幅较大，原因在于 $R^2$ 以方差为分母，对测试折内部的均值差异更敏感，而 RMSE 与 MAE 衡量的是绝对误差，二者结合更能反映实际预测水平。交叉验证 $R^2$ 在 50 个测试折上的标准差为 0.1008，说明不同折之间的表现差异明显，单条帖子的曝光量受随机传播影响很大。MAPE 未报告，是因为曝光量从 570 次到 1,110,282 次跨越三个数量级，少数极端帖子的相对误差会主导平均值，指标失去可比性。因此，模型适合解释平均效应与做情景比较，不适合对单条帖子给出精确的点预测。

残差诊断的数值已在表\ref{tab:q3_diag}与图\ref{fig:q3_resid}中给出：Jarque--Bera 统计量为 120.36、$p<0.001$，残差偏离正态；Breusch--Pagan 统计量为 54.74、$p=0.0005$，White 统计量为 266.52、$p=0.030$，两项都拒绝同方差，这正是采用 HC3 标准误的理由。Ramsey RESET 的 $F$ 值为 2.845、$p=0.059$，在 5\% 水平未检出函数形式误设；Durbin--Watson 为 2.052，接近 2；Cook 距离最大值为 0.0377，远小于 1，杠杆值超过 $2p/n=0.104$ 的观测有 10 个，属于高杠杆但未构成强影响。由此可见，残差的非正态与异方差主要影响点预测的精度，系数与区间在 HC3 标准误下仍然可用。

\subsection{设定敏感性}

灵敏度分析通过改变模型设定与缺失处理方式，观察核心效应是否发生方向反转或量级跳变。基准 M2 的付费提升为 40.13\%、Video 倍数为 4.90；加入粉丝规模（M1）后为 40.09\% 与 4.89，变化小于 0.1 个百分点；剔除月份（M3）后为 49.65\% 与 5.25，上升约 9.5 个百分点；线性月份（M4）为 47.14\% 与 5.41；Huber 回归为 37.44\% 与 5.38；Gamma 模型为 44.84\% 与 3.52。删除缺失行后付费提升为 39.98\%，与基准差 0.15 个百分点。七种设定下付费提升介于 37.44\% 至 49.65\%，全部为正且显著；视频倍数介于 3.52 至 5.41，全部大于 1 且显著。

对结果影响最大的是月份控制：剔除月份后付费提升上升约 9.5 个百分点，说明月份固定效应吸收了一部分与付费相关的季节差异，若不控制月份会把季节效应误记为付费效应。加入粉丝规模几乎不改变任何结果，再次说明该变量在月份固定效应之外没有独立信息。更换估计方法带来的变化不超过 5 个百分点，缺失处理带来的变化仅 0.15 个百分点。整体看，核心结论对模型设定不敏感，但对是否控制时间维度中等敏感，这也印证了选择月份固定效应的必要性。

\subsection{综合判断}

结合上述检验，最终模型在样本外的绝对误差仅比样本内高约 5\%，核心效应在七种扰动下方向一致、量级稳定，说明结论具备可重复性。需要保留的两点限制是：视频类型只有 7 条样本，其倍数区间很宽，精确幅度不宜过度解读；交叉验证 $R^2$ 仅约 0.10，说明曝光量的绝大部分变异来自模型之外的推荐机制与传播随机性，本报告的模型应定位为平均效应解释工具，而不是精确预测工具。
